\documentclass[10pt,a4paper]{amsart}
\usepackage[margin=19mm]{geometry}
\usepackage{amsmath,amssymb,mathtools}
\usepackage[scr=rsfs,cal=euler]{mathalfa}
\usepackage{xfrac}
\usepackage[unicode,hidelinks,bookmarks=false]{hyperref}
\newcommand{\ie}{i.e.\ }
\newcommand{\cf}{cf.\ }
\newcommand{\resp}{respectively\ }
\newcommand{\Subsection}{\subsection}
\providecommand{\nobreakdash}{\nobreak\hbox{-}}
\setlength{\emergencystretch}{2em}
\multlinegap0pt
\usepackage{amssymb}
\newtheorem{proposition}{Proposition}
\title{Isolated Singularities and Measure Data Problems for Semilinear Equations Driven by Stable L\'evy Operators}
\author{Kamil Dunst, Tomasz Klimsiak}
\date{Source: arXiv:2607.27954v1 (CC BY 4.0)}
\begin{document}
\maketitle

\noindent\emph{Source and licence.} Isolated Singularities and Measure Data Problems for Semilinear Equations Driven by Stable L\'evy Operators. Version arXiv:2607.27954v1 (CC BY 4.0), distributed by arXiv under the Creative Commons Attribution 4.0 International licence, http://creativecommons.org/licenses/by/4.0/, as recorded in the arXiv licence metadata for this exact version and shown on the abstract page https://arxiv.org/abs/2607.27954v1. Full licence notice: \texttt{licenses/arxiv-2607.27954-CC-BY-4.0.txt}.\par\smallskip

\noindent\textit{Editorial scope.} Counted unit: the article's proposition prop.spkg (source0.tex lines 3017--3024). The article's complete original proof of it is delivered with it (source0.tex lines 3026--3155, one \texttt{proof} environment of 130 lines). No mathematics is written, repaired or re-derived: the statement and the proof are the article's own lines, in the article's own notation, and no retained line is substituted or re-flowed.  No further source-local context is needed: every cross-reference of every retained line resolves inside the two delivered spans.  Nothing was omitted from any retained span, so the package carries no omission record.  No retained line contains a citation, so the wrapper carries no bibliography and no bibliography entry.  No retained line contains a figure environment, an image-inclusion command or drawing code, so no image is delivered and the package declares no asset dependency.  Editorial formatting only, in addition: an AMS \texttt{amsart} preamble that restates the article's own declarations and macros used by the retained lines, namely the environments \texttt{proposition} on separate counters, together with the standard packages those lines need.  The retained environments print in this document as Proposition 1 in order of appearance, whereas the article prints the same items with the numbering of its own full document.  The complete native source of the article as distributed in the arXiv e-print for this exact version is bundled unchanged as \texttt{sources/206376/source0.tex}.
\begin{proposition}
\label{prop.spkg}
For every $r>0$,
\[
G_{B(0,r)}(x,z)>0,
\qquad x,z\in B(0,r).
\]
\end{proposition}

\begin{proof}
Let $x,z\in B(0,r)$, $x\neq z$ and $h>0$ be such that  $B(z,h)\subset\subset B(0,r)\setminus \{x\}$. Set 
$e:=\frac{z-x}{|z-x|}$ and  $K:=[x,z]$.
Set
\[
\rho:=\frac14[(r-|z|)\wedge (r-|x|)\wedge h].
\]
Choose $\varepsilon \in (0,1/2)$ such that $|z-x|\varepsilon<\rho$.
By the strict positivity of $G(\cdot,\cdot)$,
\[
\int_0^\infty
\mathbb{P}_0\bigl(X_t\in B(e,\varepsilon)\bigr)\,dt
=
\int_{B(e,\varepsilon)}G(w,0)\,dw>0.
\]
Consequently, there exists $u_0>0$ such that
\begin{equation}
\label{eq.endpoint-positive}
\mathbb{P}_0\bigl(X_{u_0}\in B(e,\varepsilon)\bigr)>0.
\end{equation}
Since the paths of $X$ are c\`adl\`ag, they are bounded on
$[0,u_0]$ almost surely. Therefore,
there exists $M>0$ such
that
\begin{equation}
\label{eq.basic-event}
p:=
\mathbb{P}_0\left(
X_{u_0}\in B(e,\varepsilon),\
\sup_{0\leq t\leq u_0}|X_t|\leq M
\right)>0.
\end{equation}
Choose $n\in\mathbb{N}$ so large that $\frac{|z-x|M}{n}<\rho$,
and define
\[
\lambda:=\frac{|z-x|}{n},\quad
T:=\lambda^{2s} u_0,
\quad
t_k:=kT,
\quad k=0,\ldots,n.
\]
By strict $2s$-stability,
\[
\bigl(\lambda^{-1}X_{\lambda^{2s} t}\bigr)_{t\geq0}
\overset{d}{=}
(X_t)_{t\geq0}.
\]
Hence \eqref{eq.basic-event} implies
\begin{equation}
\label{eq.scaled-event}
\mathbb{P}_0\left(
X_{T}\in B(\lambda e,\lambda\varepsilon),\
\sup_{0\leq t\leq T}|X_t|\leq\lambda M
\right)=p.
\end{equation}
For $k=1,\ldots,n$, define
\[
A_k:=
\left\{
X_{t_k}-X_{t_{k-1}}
   \in B(\lambda e,\lambda\varepsilon)
\right\}
\cap
\left\{
\sup_{t_{k-1}\leq t\leq t_k}
|X_t-X_{t_{k-1}}|
\leq\lambda M
\right\},
\]
and put
\[
E_k:=\bigcap_{j=1}^{k}A_j,
\qquad E_0:=\Omega.
\]
By independence and stationarity of increments of the L\'evy process
\begin{equation}
\label{eq.iterated-probability}
\mathbb{P}_x(E_k)=p^k,
\qquad k=1,\ldots,n.
\end{equation}
We next verify that on $E_n$ the process stays inside $B(0,r)$ up to
time $t_n$ and reaches $B(z,h)$ at time $t_n$. For $k=1,\ldots,n$ and  $t\in[t_{k-1},t_{k}]$.
\[
X_{t}-X_0
=
X_t-X_{t_{k-1}}+\sum_{j=1}^{k-1}
\bigl(X_{t_j}-X_{t_{j-1}}\bigr).
\]
Therefore,  on $E_n$, using the definition of $A_k$, we obtain, under the measure $\mathbb P_x$,
\[
\begin{split}
\left|
X_t-[x+(k-1)\lambda e]
\right|&=\left|
X_t-X_0-(k-1)\lambda e
\right|
\\&\leq
|X_t-X_{t_{k-1}}|+
\left|
X_{t_{k-1}}-X_0
-(k-1)\lambda e
\right|\\&\leq
\lambda M+(k-1)\lambda\varepsilon<2\rho.
\end{split}
\]
Thus, by the choice of $\rho$ and $\lambda$,
\[
X_t\in B(0,r),
\quad 0\leq t\leq t_n,
\]
and
\[
|X_{t_n}-z|=|X_{t_n}-[x+n\lambda e]|=|X_{t_n}-X_0-n\lambda e| \le  n\lambda\varepsilon
=
|z-x|\varepsilon
<\rho
\]
on $E_n$, which in turn implies that
\[
X_{t_n}\in B(z,h),\quad t_n<\tau_{B(0,r)}
\]
on $E_n$. Consequently, by right continuity of the paths
\[
\mathbb{E}_x
\int_0^{\tau_{B(0,r)}}\mathbf{1}_{B(z,h)}(X_t)\,dt=\int_{B(z,h)} G_{B(0,r)}(x,y)\,dy>0.
\]
Since this holds  for any $z\in B(0,r)\setminus \{x\}$ and $h>0$ such that $B(z,h)\subset\subset B(0,r)\setminus\{x\}$,
we conclude that $G_{B(0,r)}(x,\cdot)$ is strictly positive a.e. in $B(0,r)$. This combined with the fact that $G_{B(0,r)}(x,\cdot)$
is $(P^{*,B(0,r)}_t)$-excessive  yields the desired result.  
\end{proof}

\end{document}
